\section{Preliminaries}\label{sec:prelims-sos} 

\paragraph{Notations.} 
For a (weighted) graph $G=(V,E)$, we denote by $(u,v) \sim E$ the distribution over edges proportional to their weight, and by $A_G$ the transition matrix related to the random walk on $G$.
We define the Laplacian of $G$ by $L_G = I - A_G$, and denote the stationary distribution of $G$ by $\pi_G$. In this paper we will mostly be working with regular graphs and therefore the uniform distribution over $G$. We use $\ip{f,g}_G$ to denote $\E_{u \sim \pi_G}[f(u)g(u)]$. We often drop the subscript of $G$ when it is clear from context. 
%\footnote{By that, we mean the distribution that assigns a vertex $v$ probability distribution which is proportional to $\sum\limits_{u}w(u,v)$.} 

%It is easy to see that picking a random edge from $E$ is equivalent to picking a random vertex $v \sim \pi$, and then a random neighbor $w$ of $v$ with probability proportional to the weight of the edge $(w,v)$. For $v \in V(G)$, we denote $\deg_G(v) = \sumlimits_{u} w(u,v)$.
If $A$ is some probabilistic event or condition, we use $\Ind(A)$ to denote the indicator random variable of $A$ (i.e., $\Ind(A)=1$ if $A$ occurs and $\Ind(A)=0$ otherwise). 


% \begin{definition}[Unique games] \label{def:uniguegames}
% A \emph{unique games instance} is a pair $I=(G,\Pi)$ where $G=(V,E)$ is a graph and $\Pi$ is a collection $\{ \pi_{u,v} \}_{(u,v) \in E}$ such that $\pi_{u,v}$ is a permutation over some finite set $\Sigma$. The graph $G$ is known as the \emph{constraint graph} of $I$.

% Given an instance $I=(G,\Pi)$ of unique games and an assignment $x\in \Sigma^V$ of values to the vertices of $G=(V,E)$, the   \emph{value} of $x$ with respect to $I$ is $\val_I(x) = \E_{(u,v) \sim E} \Ind(\pi_{u,v}(x_u) = x_v)$.
%  The value of $I$ is the maximum of $\val_I(x)$ over all $x\in \Sigma^V$.
% We may drop the subscript $I$ when the instance is clear from context.

% We say that $(G,\Pi)$ is an \emph{affine} unique games instance if $\Sigma$ is an additive group and all the functions $\pi_{u,v}$ are of the form $\pi_{u,v}(x) = x -  a_{u,v}$ for some $a_{u,v} \in \Sigma$. That is, all constraints correspond to $x_u - x_v = a_{u,v}$. 
% \end{definition}

% It is known that the UGC is equivalent to its restriction on affine instances~\cite{KhotKMO07}. 
% In this paper we restrict attention to affine instances only.
% For the sake of simplicity, we will drop the qualifier ``affine'' in future discussion, but all of our results are for this family of constraints. 

\subsection{The Sum of Squares Relaxation of Unique-Games}
Our algorithm is based on the \sos semidefinite programming (SDP) relaxation, and in particular its view as optimizing over \emph{pseudo expectation} operators. In this section, we 
briefly present the necessary background as well as several 
basic facts that we will use, and we refer the reader to the
surveys \cite{BarakS14,RaghavendraSS18,TCS-086} for a more 
systematic presentation.

Given a Unique-Games instance $I= (G=(V,E),\Sigma,\Pi)$
the value of $I$ can be computed by the following integer program over zero-one variables $\{ X_{u,a} \}_{u \in V, a \in \Sigma}$:
\begin{align}
\max_{X}& \E_{(u,v) \in E} \sum_{a \in \Sigma} X_{u,a} X_{v,\pi_{uv}(a)} \label{eq:ip}\\
s.t. 
&\quad X_{u,a}^2 = X_{u,a} ~\qquad\quad \forall u \in V, a \in \Sigma \nonumber,\\
&\quad X_{u,a}X_{u,b} = 0 \qquad\quad \forall u \in V, a \neq b \in \Sigma \nonumber,\\
&\quad \sum_{a} X_{u,a} = 1 ~\qquad\quad \forall u \in V\nonumber.
\end{align}

Indeed, an assignment to the variables $X_{i,a}$ represent 
an assignment to the vertices of $G$, in which a vertex $v \in V$ takes label $a \in \Sigma$ if $X_{v,a} = 1$, and the objective function counts the fraction of constraints 
that are satisfied satisfied.


The degree $D$ \sos relaxation of the above program is a convex 
optimization program that relaxes the above program. This 
program can be written as an optimization problem in which 
the goal is to find a vector-valued assignment to the variables $X_{S,\sigma}$, 
where we have such variable for each set of vertices $S$ of size at most $D$ and labels for them $\sigma\in \Sigma^S$. 
The constraints of the program address the inner products between these vectors, attempting to ensure that for any two sets 
of vertices $S$ and $T$ such that $|S\cup T|\leq D$, 
the inner products 
$\langle X_{S,\sigma}, X_{T,\sigma'}\rangle$ represent 
a distribution over assignments to $S\cup T$, and that 
these distributions are locally consistent. 
It will be more convenient for us to take a different 
but equivalent view on the solutions to the \sos relaxation,
in the language of pseudoexpectations as presented next.

\subsection{Pseudoexpectations, Pseudodistributions and Pseudoprobabilities}
\paragraph{Pseudoexpectations.}
A degree $D$ pseudoexpectation $\pE$ is an operator $\pE:X^{\le D} \to \R$, where $X^{\le D}$ is the set of all monomials in the $X$ variables up to degree $D$, and $\pE$ satisfies the above equality constraints and the Booleanity constraints $\{X_{u,a}^2 = X_{u,a}\}$ as axioms.  For brevity, we will refer to this set of axioms as $\A_{I}$, dropping the subscript when $I$ is clear from context. More generally, given a polynomial optimization program $P = \{\max_{x} p(x) \ s.t. ~ q_i(x) = 0, \forall i \in [m] \}$, the degree-$D$ sum-of-squares semidefinite programming relaxation of $P$ is a semidefinite program of size $n^{O(D)}$ that returns a {\em pseudoexpectation operator} $\pE: x^{\le D} \to \R$. This operator can be uniquely extended to give a pseudoexpectation operator on the set of all polynomials of degree at most $D$ by linearity (defined precisely below).
This operator satisfies four properties:
\begin{itemize}
\item Scaling: $\pE[1] = 1$.
\item Linearity: $\pE[a \cdot f(x) + b \cdot g(x)] = a\cdot \pE[f(x)] + b \cdot \pE[g(x)]$, for all $a,b \in \R$ and all degree $\leq D$ polynomials $f,g$. 
\item Non-negativity of low-degree squares: $\pE[s(x)^2] \ge 0$ for all polynomials $s(x)$ with $\deg(s) \le \tfrac{D}{2}$.
\item Program constraints: $\pE[f(x) \cdot q_i(x)] = 0$ for all $i \in [m]$ and polynomials $f(x)$ such that $\deg(fq_i) \le D$.
\end{itemize}
Depending on the problem we are trying to solve, we also 
discuss the value achieves by the pseudoexpectation $\pE$, 
which is defined to be the pseudoexpectation of the objective 
function. In our case of interest, namely the case of Affine 
Unique-Games, the \emph{value} of $\pE$ on the instance $I$ is denoted by $\val_\mu(I)$ and is defined to be 
\[
\pE[\val_{I}(X)] = \pE \left[ \E_{(u,v) \in E} \sum_{a \in \Sigma} X_{u,a} X_{v,\pi_{uv}(a)}\right].
\]
Our pseudoexpectation will be guaranteed to achieve a value which matches the completeness guarantee (for example, in 
the context of Theorems~\ref{thm:main_johnson},~\ref{thm:main_grass} it will satisfy that $\pE[\val_{I}(X)]\geq c$).

\paragraph{Pseudodistribution.} 
We often refer to a pseudoexpectation operator $\pE$ as an operator corresponding to a pseudodistribution $\mu$ over variables $X$. This is analogous to the case where we have an actual distribution $\mu$ and its corresponding expectation operator $\E_\mu[\cdot]$. This notation makes our analysis using $\pE$ operators more intuitive, since many properties that are true of actual distributions also hold for pseudodistributions. Hence when we say that we are given a degree $D$ pseudodistribution $\mu$ we are referring to the degree $D$ pseudoexpectation operator $\pE_\mu$.

% A pseudoexpectation operator $\pE$ 
% has an associated real-valued pseudodistribution $\mu$ so that 
% \[
% \pE[ f ]= \pE_{\mu}[ f ] = \sum\limits_{\omega} \mu(\omega) f(\omega). 
% \]


\paragraph{Pseudoprobabilities.}
\begin{definition}[Pseudoprobability of an event]\label{def:pseudoprob}
Let $\mu$ be a pseudodistribution of degree $D$.
For an event $\cE(X,X')$ such that $\Ind[\cE(X)]$ can be expressed as a degree-$D$ function of $X$, we define the {\em pseudoprobability of $\cE(X)$} to be
\[
\pPr_{\mu}[\cE(X)] = \pE_{\mu}[\Ind(\cE(X)].
\]
Similarly, if $\cF(X)$ is an event and $\deg(\Ind[\cF(X)]) + \deg(\Ind[\cE(X)]) \le D$,  we define the {\em pseudoprobability of $\cE(X)$ conditioned on $\cF(X)$} to be
\[
\pPr_\mu[\cE(X) \mid \cF(X)] = \pE_\mu[\Ind(\cE(X)) \mid \Ind(\cF(X))] = \frac{\pE[\Ind(\cE(X)) \cdot \Ind(\cF(X))]}{\pE[\Ind(\cF(X))]}.
\]
\end{definition}

% More generally, given a polynomial optimization program $P = \{\max_{x} p(x) \ s.t. ~ q_i(x) = 0, \forall i \in [m] \}$, the degree-$D$ sum-of-squares semidefinite programming relaxation of $P$ is a semidefinite program of size $n^{O(D)}$ that returns a {\em pseudoexpectation operator} $\pE: x^{\le D} \to \R$. This operator can be uniquely extended to give a pseudoexpectation operator on the set of all polynomials of degree at most $D$ by linearity (defined precisely below).
% This operator satisfies four properties:
% \begin{itemize}
% \item Scaling: $\pE[1] = 1$.
% \item Linearity: $\pE[a \cdot f(x) + b \cdot g(x)] = a\cdot \pE[f(x)] + b \cdot \pE[g(x)]$, for all $a,b \in \R$ and all degree $\leq D$ polynomials $f,g$. 
% \item Non-negativity of low-degree squares: $\pE[s(x)^2] \ge 0$ for all polynomials $s(x)$ with $\deg(s) \le \tfrac{D}{2}$.
% \item Program constraints: $\pE[f(x) \cdot q_i(x)] = 0$ for all $i \in [m]$ and polynomials $f(x)$ such that $\deg(fq_i) \le D$.
% \end{itemize}
% Additionally, we will have $\pE[p(x)] \ge {\sf vap}(P)$.

% The \emph{value} of $\pE$ on the instance ${I}$, 
% denoted by $\val_\mu({I})$, is defined to be 
% \[
% \pE[\val_{I}(X)] = \pE \left[ \E_{(u,v) \in E} \sum_{a \in \Sigma} X_{u,a} X_{v,\pi_{uv}(a)}\right].
% \]

% \subsection{Sum of Squares}
% Given a polynomial optimization program $P = \{\max_{x} p(x) \ s.t. ~ q_i(x) = 0, \forall i \in [m] \}$, the degree-$D$ sum-of-squares semidefinite programming relaxation of $P$ is a semidefinite program of size $n^{O(D)}$ that returns a {\em pseudoexpectation operator} $\pE: x^{\le D} \to \R$. This operator can be uniquely extended to give a pseudoexpectation operator on the set of all polynomials of degree at most $D$ by linearity (defined precisely below).
% This operator satisfies four properties:
% \begin{itemize}
% \item Scaling: $\pE[1] = 1$.
% \item Linearity: $\pE[a \cdot f(x) + b \cdot g(x)] = a\cdot \pE[f(x)] + b \cdot \pE[g(x)]$, for all $a,b \in \R$ and all degree $\leq D$ polynomials $f,g$. 
% \item Non-negativity of low-degree squares: $\pE[s(x)^2] \ge 0$ for all polynomials $s(x)$ with $\deg(s) \le \tfrac{D}{2}$.
% \item Program constraints: $\pE[f(x) \cdot q_i(x)] = 0$ for all $i \in [m]$ and polynomials $f(x)$ such that $\deg(fq_i) \le D$.
% \end{itemize}
% Additionally, we will have $\pE[p(x)] \ge value(P)$.
% We refer to these as {\em pseudomoments} of a {\em pseudodistribution}.
% \dor{basically copy-pasted above; can we remove this?}


\subsection{SoS-ing Mathematical Proofs}
Our argument will use several mathematical statements (such 
as global hypercontractivity), and we will need to be 
able to argue that these statements are also satisfied 
in the context of the SoS program and its variables. Thus, 
we will need to be able to ensure that we use tools that can 
be proved via sum of squares inequalities of low-degree (and 
axioms), and below we collect a few such standard tools that we will use. We use the following notation for SoS prooofs.

\paragraph{Sum of squares proofs notations.}
Given a set of axioms $\cA = \{q_i = 0\}_i \cup \{g_j \ge 0\}_j$ for polynomials $q_i, g_j \in \R[X]$, we say that ``there is a degree-$d$ sum-of-squares proof that $f \ge h$ modulo $\cA$'' if: $f = h + s + \sum_k c_k\cdot Q_k + \sum_t r_t \cdot G_t$ where each polynomial $Q_k$ and $G_t$ is a product of some polynomials from $\{q_i\}$ and $\{g_j\}_j$ respectively, $s,\{c_k\}_k,\{r_t\}_t \in \R[X]$ are real polynomials such that $s$ and $\{r_t\}_t$ are sums of squares, and the maximum degree among $s,\{c_k Q_k\}_k, \{r_t G_t\}_t$ is at most $d$. We will use the notation $\cA \vdash_d f(x) \ge h(x)$ to denote the existence of such an equality.
We also sometimes use $f(x) \succeq h(x)$ to denote that the inequality is a \sos inequality.


\subsubsection{Basic Inequalities}
The first of which is the following basic forms of the Cauchy-Schwarz and H\"{o}lder inequalities. 
The proofs are by now standard and can be found for example in~\cite{BarakKS14,ODonnellZ13}.
\begin{lemma}[Cauchy Schwarz]\label{CS1-prelim}
For for all $\epsilon \in \R_+$, 
$$
\vdash_2 YZ \leq \frac{\epsilon}{2}Y^2 + \frac{1}{2\epsilon}Z^2.
$$
\end{lemma}

\begin{lemma}[Cauchy Schwarz]\label{CS2-prelim}
A degree-$D$ pseudoexpectation operator where 
$D \geq  2\max(\deg(f),\deg(g))$ satisfies that
$$\pE[fg]^2 \leq \pE[f^2]\pE[g^2].$$
\end{lemma}
% Proofs for both lemmas appear in \cite{BarakKS14}.


% We also need the following version of H\"{o}lder's inequality which is proven in e.g.~\cite{ODonnellZ13}.
\begin{fact}[H\"{o}lder's Inequality]\label{fact:sos-hol}
For all real $\nu > 0$ we have that,
$$\vdash_4 Y^3 Z \leq \frac{3\nu}{4} Y^4 + \frac{1}{4\nu^3}Z^4.$$ 
\end{fact}



The theory of univariate sum-of-squares (in particular, Luk\'{a}cs Theorem) says that if a univariate polynomial is non-negative on an interval, this fact is also SoS-certifiable. The following  corollary of Luk\'{a}cs theorem is well-known.
\begin{corollary}[Corollary of Luk\'{a}cs Theorem]\label{cor:Lukacs}
Let $q$ be a degree-$d$ polynomial which is non-negative on $[a,b]$.
Then given the axioms $\cA = \{x \ge a\} \cup\{x \le b\}$, there is a degree-$2d$ \sos proof that $q$ is non-negative, $\cA \vdash_{2d} q(x) \ge 0$.
\end{corollary}
We will use the above multiple times to convert univariate inequalities into \sos inequalities in a blackbox manner.


Given multivariate inequalities the theorem above no longer holds. Nevertheless given a strictly positive polynomial $f$ that is bounded away from $0$, one can get an \sos proof of degree which is exponential in $\deg(f)$ using Theorem 3 in~\cite{schweighofer}. We state a corollary of~\cite{schweighofer} that is sufficient for our purposes. 

\begin{theorem}[Corollary of~\cite{schweighofer}]\label{thm:blackbox-sos}
Let $S$ be the set $[0,1]^k$. Let $f(x)$ be a polynomial of degree $d$ with $\nu = \min\{f(x) \mid x \in S\} > 0$ and $||f||$ denoting the maximum absolute value of $f$'s coefficients. Then $f$ has an \sos certificate of bounded degree:
\[\{x_i \in [0,1] | i \in [k]\} \vdash_D f(x) \geq 0,\]
with $D = \poly(k^{\deg(f)}, ||f||/\nu)$.
\end{theorem}

We will use the above fact in Section~\ref{sec:johnson_2} to convert inequalities involving polynomials of constant degree and on constantly many variables into \sos inequalities. We believe it should be possible to get a degree bound above which is $\poly(\deg(f))$ in cases where we have reasonable polynomials $f$, though as far as we know this has not been proved in generality.




% The next claim asserts that sum of squares proofs can capture 
% inequalities related to projection operators onto eigenvalue 
% spaces of transition operators.
% \begin{claim}\label{claim:spectral-sos}
% Let $A$ be the transition matrix for a random walk on an undirected (weighted) graph $G$ and $\pi$ be the stationary measure on $G$. Then $A$ has real eigenvalues. 
% Moreover, letting $\Pi$ be the projection operator onto 
% the space of $A$'s spanned by eigenvectors with eigenvalues 
% at most $\lambda$, there is a degree-$2$ \sos proof that
% \[
% \langle f, A \Pi f \rangle_{\pi} \leq \lambda \langle f, \Pi f \rangle_{\pi}
% \]
% and 
% \[
% \langle f, A (\Id - \Pi) f \rangle_{\pi} \leq \langle f, (\Id - \Pi) f\rangle_{\pi}.
% \]
% \end{claim}
% \begin{proof}
% We use that $A$ is self-adjoint in the inner product space $(\R^n,\langle \cdot \rangle_\pi)$, and therefore it has real eigenvalues and its right eigenspace has orthonormal eigenvectors $v_1,\ldots,v_n$.
% We may write $f$ according to its orthogonal decomposition, $f = \sum_{i = 1}^n c_i \cdot v_i$ for $c_i$ linear functions of $f$, and if there are $k$ eigenvalues of value at most $\lambda$ then $\Pi f = \sum_{i = 1}^k c_i \cdot v_i$.
% We thus have
% \begin{align*}
% \langle f, A\Pi f \rangle_\pi
% &= \left\langle \sum_{i = 1}^n c_i v_i, \sum_{j=1}^k c_j \cdot A v_j\right\rangle_\pi\\
% &= \left\langle \sum_{i = 1}^n c_i v_i, \sum_{j=1}^k c_j \lambda_j \cdot v_j\right\rangle_\pi
% = \sum_{j = 1}^k \lambda_j c_j^2
% \le \lambda \sum_{j=1}^k c_j^2
% = \lambda \langle f, \Pi f \rangle_\pi,
% \end{align*}
% where the inequality is a degree-2 sum of squares because $\lambda_j \le \lambda$, and the $c_j$ are degree-1 functions of $f$.
% A near-identical proof gives the second statement when we observe that $A$'s maximum eigenvalue is $\le 1$.
% \end{proof}



\subsubsection{Approximating Indicators via Low-degree Polynomials}\label{sec:apx-ind}
% In this section, we note that there is a low-degree polynomial which provides an SOS-certifiably good approximation to a step function.
% This will be a consequence of the existence of low-degree approximations to step functions that appear in the literature, as well as the theory of univariate sums-of-squares.
%As for the \sos-certifiability, this follows as a simple corollary from Luk\'{a}cs' Theorem as we explain next.


% \begin{theorem}[Luk\'{a}cs Theorem, see e.g. \cite{szeg1939orthogonal}]
% If $p$ is a degree-$d$ univariate polynomial with $p(x) \ge 0$ for $x \in [-1,1]$, then $p$ can be written as 
% \[
% p(x) = \begin{cases}
% s(x)^2 + (1-x^2) t(x)^2 & \text{ if $d$ even}\\
% (1+x)s(x)^2 + (1-x) t(x)^2 & \text{ if $d$ odd}
% \end{cases}
% \]
% for $s,t$ real polynomials of degree at most $d$.
% \end{theorem}

Our argument will involve indicators of events such as 
$f(X)\geq \beta$ where $f$ is a low-degree polynomial, and we will want to condition on such events. Strictly speaking, the function $\Ind[f(X)\geq \beta]$ is not a low-degree polynomial and therefore we cannot condition on it. However, it is not difficult to show that such indicators can be approximated by low-degree polynomials, and we will need to use such ideas. 
Indeed, in this section we present such an approximation theorem that will be used throughout our proofs.

The following theorem, due to~\cite{diakonikolas2010bounded}, provides a low-degree approximation to a step function.
%Similar statements may be proven using classical results in approximation theory, however the degree bounds in~\cite{diakonikolas2010bounded} are sharper.


\begin{theorem}\label{thm:step-approx}
Let $\Ind[x \geq \beta]$ be the step function at $\beta \in (0,1)$. Then for each $0 < \nu < \beta$ there is a univariate polynomial $p_{\beta,\nu}$ of degree $O(\frac{1}{\nu}\log^2\frac{1}{\nu})$ such that:
\begin{enumerate}
\item $|p_{\beta,\nu}(x) - \Ind[x \geq \beta]| \le \nu$ for all $x \in [0, \beta] \cup [\beta +\nu, 1]$.
\item $p_{\beta,\nu}$ is monotonically increasing on $(\beta, \beta + \nu)$.
\item $0 \le p_{\beta,\nu}(x) \le 1$ for all $x \in [0,1]$.
\item All coefficients of $p_{\beta,\nu}$ are at most
$2^{O(\log(1/\nu)/\nu)}$ in absolute value.
\end{enumerate}
Further the first three facts are SoS-certifiable in degree $2\deg(p_{\beta,\nu})$.
% given axioms $\cA = \{x \ge 0\}\cup \{x \le 1\}$, there is an SoS proof that 
% \[ A \vdash_{O(\frac{1}{\delta}\log^2 \frac{1}{\epsilon})} \{0 \le p_{\alpha}^{\epsilon,\delta}(x) \le 1\}.
% \]
\end{theorem}
The first three items follow from~\cite[Theorem 4.5]{diakonikolas2010bounded}, the fourth 
item is immediate by Markov brothers' inequality, and the SoS certifiability follows from Corollary~\ref{cor:Lukacs}.

The following fact provides convenient point-wise bounds on the approximating polynomials of indicators from above. %Theorem~\ref{thm:step-approx}.
\begin{fact}[Markov Inequality for Bounded Polynomials]
\label{fact:bdd-markov}
Let $p := p_{\beta,\nu}$ be the degree-$D = \tO(1/\nu)$ polynomial guaranteed by Theorem \ref{thm:step-approx}.
Then $p$ satisfies Markov's inequality:
\begin{equation*}
\{0 \le x \le 1\}
\vdash_{\deg(p)} 
\{ p(x) \ge 1 - \frac{1-x}{1-\beta-\nu} - \nu\} \cup \{p(x) \le \frac{x}{\beta-\nu} + \nu\}
\end{equation*}
\end{fact}
\begin{proof}
We perform case analysis on $x$ and then use Corollary~\ref{cor:Lukacs} to conclude the proof
is \sos.
For the first inequality, for $x \in [0,\beta+\nu)$ 
we have
\[
p(x) \succeq 0 \succeq 1 - \frac{1-x}{1-\beta - \nu},
\]
where we have used that $p(x) \succeq 0$ and $\frac{1-x}{1-\beta - \nu}\succeq 1$. For $x \in [\beta +\nu, 1]$,
\[
p(x) \succeq 1 - \eps \succeq 1 - \eps - \frac{1-x}{1-\alpha-\delta},
\]
where we have used that $x \in [0,1]$ so that we are subtracting a positive quantity.
Combining these claims concludes the proof of the first claim.

To see the second claim, notice that for $x \in [0,\beta-\nu]$, $p(x) \le \nu$, and for $x \in (\beta - \nu, 1]$, $p(x) \le 1 \le \frac{x}{\beta - \nu}$.
This concludes the proof.
\end{proof}

% We remark that the above statement is very similar but not identical to~\cite[Theorem 4.5]{diakonikolas2010bounded}.
% To switch from their $\text{sign}(y)$ polynomial for $y \in [-1,1]$ to $s_\alpha(x)$ for $x \in [0,1]$, we can do a simple change of variables, taking $y = x-\alpha$.
% Shifting by a constant and rescaling changes the bounds so that $p(x) \in [0,1]$.
% The third item is not explicitly written in the statement of Theorem 4.5 of \cite{diakonikolas2010bounded}, but it can be easily extracted from the proof therein. 

% \begin{proof}
% We claim that Luk\'{a}cs theorem implies
% \[
% q(x) = \begin{cases}
% s(x)^2 + (x-a)(b-x)t(x)^2 & \text{ if $d$ even}\\
% (x-a)s(x)^2 + (b-x)t(x)^2 & \text{ if $d$ odd}
% \end{cases}
% \] 
% for $s,t$ real polynomials of degree at most $2d$, and this implies our corollary.
% To get the claim, we perform a change of variables, taking $x' =\frac{2}{b-a} (x - a) -1$.
% Let $q(x) = h(x')$. 
% We now have that $h(x')$ is a degree-$d$ polynomial which is non-negative on $[-1,1]$.
% From Luk\'{a}cs Theorem, we have that
% \[
% h(x') = \begin{cases}
% u(x')^2 + (1-x'^2) v(x')^2 & \text{ if $d$ even}\\
% (1+x')u(x')^2 + (1-x') v(x')^2 & \text{ if $d$ odd}
% \end{cases}
% \]
% for $u,v$ real polynomials of degree at most $d$.
% But now,
% \[
% \tfrac{b-a}{2} \cdot (1 - x') = b-x, \qquad \tfrac{b-a}{2}\cdot (1 + x') = x-a, \qquad \left(\tfrac{b-a}{2}\right)^2\cdot (1-x'^2) 
% = (x-a)(b-x)
% \]
% and so by applying a change of variables to the polynomials $u(x'),v(x')$ to obtain $s(x),t(x)$, the conclusion follows.
% \end{proof}


% When we have \sos certificates that polynomials are bounded within $(0,1) \pm \epsilon$, \sos can also certify that they behave roughly like indicator functions.
% \begin{fact}[Union bound for Approximate Indicators]\label{fact:apx-ub}
% Suppose that $h,g$ are polynomials of degree at most $d$, and suppose furthermore that from the axioms $A$, there is an \sos proof that $A \vdash_d \{0 \le g \le 1\}\cup\{ 0 \le h \le 1\}$.
% Then,
% \[
% A \quad \vdash_{2d}\quad gh \ge g + h - 1
% \]
% \end{fact}
% \begin{proof}
% We have as a polynomial equality that $(g + (1-g))(h + (1-h)) = 1$.
% Expanding then re-arranging,
% \begin{align*}
% gh 
% &= 1 - g(1-h) - h(1-g) - (1-g)(1-h)\\
% &\succeq 1 - g(1-h) - h(1-g) - (1-g)(1-h) - (1-h)(1-g)\\
% &= 1 - (1-h) - (1-g) 
% \end{align*}
% where in the second line we have used the \sos bounds $g,h \preceq 1$.
% Simplifying gives the conclusion.
% \end{proof}

% \begin{fact}[Markov Inequality for Bounded Polynomials]
% \label{fact:bdd-markov}
% Let $p := p_{\alpha}^{\varepsilon,\delta}$ be the degree-$D = O(\frac{1}{\delta}\log^2\frac{1}{\varepsilon})$ polynomial guaranteed by Theorem \ref{thm:step-approx}.
% Then $p$ satisfies Markov's inequality:
% \begin{equation*}
% \{0 \le x \le 1\}
% \vdash_{\deg(p)} 
% \{ p(x) \ge 1 - \frac{1-x}{1-\alpha-\delta} - \eps\} \cup \{p(x) \le \frac{x}{\alpha-\delta} + \eps\}
% \end{equation*}
% \end{fact}
% \begin{proof}
% We will perform case analysis on $x$, throughout using Corollary~\ref{cor:Lukacs} to obtain our \sos inequalities.
% We prove the first inequality first.
% For $x \in [0,\alpha+\delta)$,
% \[
% p(x) \succeq 0 \succeq 1 - \frac{1-x}{1-\alpha - \delta},
% \]
% where we have used that $p(x) \succeq 0$ and $\frac{1-x}{1-\alpha - \delta}\succeq 1$.
% Now for $x \in [\alpha +\delta, 1]$,
% \[
% p(x) \succeq 1 - \eps \succeq 1 - \eps - \frac{1-x}{1-\alpha-\delta},
% \]
% where we have used that $x \in [0,1]$ so that we are subtracting a positive quantity.
% Combining these claims concludes the proof of the first claim.

% To see the second claim, notice that for $x \in [0,\alpha-\delta]$, $p(x) \le \eps$, and for $x \in (\alpha - \delta, 1]$, $p(x) \le 1 \le \frac{x}{\alpha - \delta}$.
% This concludes the proof.
% \end{proof}

% \begin{observation}\label{obs:squaring}
% Let $p := p_{\alpha}^{\varepsilon,\delta}$ be the degree-$D = O(\frac{1}{\delta}\log^2\frac{1}{\varepsilon})$ polynomial guaranteed by Theorem \ref{thm:step-approx}.
% Then $p^2$ is a polynomial of degree $2D$ which enjoys the same guarantees as the polynomial $p_{\alpha}^{2\varepsilon,\delta}$.
% \end{observation}
% \begin{proof}
% The polynomial $p^2$ is bounded in $[0,\eps^2]$ on $[0,\alpha -\delta]$, inherits the monotonic increasing property on $(\alpha-\delta,\alpha+\delta)$, and is bounded by $[(1-\eps)^2,1]$ on $[\alpha + \delta,1]$.
% \end{proof}

\subsection{Manipulating pseudoexpectations}

\subsubsection{Reweighing and conditioning:} 
We will sometimes {\em reweigh} or {\em condition} our degree-$D$ pseudodistribution by a polynomial $s(x)$ where $s(x)$ is non-negative under the program axioms, i.e. $\cA \vdash_{d} s(x)$ for $d < D$. Technically, this operation amounts to defining a new pseudoexpectation operator $\pE'$  of degree $D-d$ by taking, 
\[\widetilde{E}'[x^{\alpha}] = \frac{\pE[x^{\alpha} \cdot s(x)]}{\pE[s(x)]},\] 
for every monomial $x^{\alpha}$ of degree at most $D - d$. As an example, under the unique games axioms presented in \ref{eq:ip} one can prove that a variable $X_i$ is in $[0,1]$, hence one can reweigh the pseudodistribution by $X_i$. One can show that reweighing preserves the four properties of the pseudodistribution up to degree $D - d$.
Thus, we will also refer to this operation as ``conditioning'', and denote $\pE'$ by $\pE[\cdot ~|~ s(x)]$. Often times, the polynomial $s(x)$ we will ``condition'' on 
will be a smooth approximation of some event $E$, in which case the above operation takes the interpretation of conditioning our sample from the pseudodistribution to satisfy some properties specified by the event $E$. We refer the reader to \cite{BarakRS11,BKS17} for further discussion of reweighting.


\subsubsection{Independent Samples}\label{sec:ind_samples}
Recall that a given pseudoexpectation operator $\pE: X^{\leq D} \rightarrow \R$ has the interpretation as averaging of functions $f(X)$ over a pseudodistribution $X\sim\mu$.  
We will need to be able to mimic averaging over two 
independently chosen samples $X,X'\sim \mu$,\footnote{Similar constructs have been used in the literature, see e.g. \cite{BarakKS14}.} for that 
we define the product pseudoexpectation $\pE_{X,X'}$ 
as follows: let $X^{\alpha}(X')^{\beta}$ be a monomial 
of degree at most $D$ in variables $X,X'$; we define  $\pE_{X,X'}[X^{\alpha}(X')^{\beta}] := \pE_{X}[X^{\alpha}] \cdot \pE_X[X^{\beta}]$. It is easy to check that $\pE_{X,X'}$ 
is also a pseudoexpectation operator corresponding to two 
independent samples from the pseudodistribution $\mu$; 
see Fact~\ref{fact:indep}.

Given two independent samples $X,X'\sim \mu$, we will often be
interested in the variables $\{Z_{v,s}\}_{v\in V, s\in\Sigma}$ 
that are used to define the shift-partition discussed in the introduction. Formally, we define them as
\[
Z_{u,s} = \sum_{a \in \Sigma} X_{u,a} X'_{u,a+s}.
\]
In the rest of this section we present a few facts about polynomials in independent samples,
and the pseudodistribution $\pE_{X,X'}$ on them. 
Some of these facts will be general, and some of which 
will be specific to the shift partition variables $Z_{u,s}$.

The first fact asserts that $\pE_{X,X'}$ defined above is a 
legitimate pseudodistribution that inherits all of the 
constraints that $\pE$ satisfies.
\begin{fact}\label{fact:indep}
If $\pE_{X}$ is a valid pseudodistribution of degree $D$ in variables $X$, then $\pE_{X,X'} $ is a valid pseudodistribution of degree $D$. Furthermore, if there are additional \sos inequalities that are true for $\T_X$, they also hold for $\T_{X,X'}$.
\end{fact}
% \begin{proof}
% The scaling and linearity of $\pE_{X,X'}$ are immediate from 
% the definition, and we next check that $\pE_{X,X'}$ is non
% negative on squares.
% This fact follows from the fact that the degree-$D$ pseudomoment matrix of $\pE_{X,X'}$ is a principal minor of the Kronecker square of the pseudomoment matrix of $X$, that is, of $(\pE X^{\le D})^{\otimes 2}$.
% Since $\pE_{X}$ is a valid pseudoexpectation matrix, $\pE_{X}[X^{\le D}]$ is a PSD matrix, and therefore its Kronecker square and any principal minor thereof. 
% Finally, in the standard manner any degree-$D$ square polynomial $s$ in variables $X,X'$ can be written as a quadratic form of $s$'s coefficient vector with the submatrix of the Kronecker square.
% Thus $\pE_{X,X'}$ satisfies the degree-$D$ \sos inequalities.


% Finally, to see that $\pE_{X,X'}[q_i(X) \cdot f(X,X')] = 0$ for any $q_i$ for which we have the constraint $q_i(x) = 0$ and any $f$ of degree at most $D - \deg(q_i)$, we write $f$ in the monomial basis, $f(X,X') = \sum_{\alpha, \beta} \hat f_{\alpha,\beta} \cdot X^{\alpha}(X')^\beta$, and then we have by linearity
% \[
% \pE_{X,X'}[q_i(X) f(X,X')] = \sum_{\alpha,\beta} \hat f_{\alpha,\beta} \cdot \pE[q_i(X) \cdot X^{\alpha}] \cdot \pE[(X')^{\beta}] = 0,
% \]
% since $\pE_X[ q_i(X) \cdot X^{\alpha}] = 0$.
% This concludes the proof.
% \end{proof}

Now, we prove some properties specific to the $Z$ variables. 
The following fact asserts that the shift partition variables 
$Z_{v,s}$ indeed behave like a partition, in the sense 
that if $X$ satisfy the constraints in the program~\ref{eq:ip}, 
then the variables $Z_{v,s}$ indeed define a partition.
\begin{fact}\label{fact:z-vars}
Define the shift variable $Z_{u,s} = \sum_{a \in \Sigma} X_{u,a}X'_{u,a+s}$ to be the indicator that $X_u - X_u' = s$, for $X,X'$ degree-$8$ solutions to the \sos relaxation of the UG integer program (\ref{eq:ip}), and for each edge $(u,v)$ the variable $Y_{(u,v)} = \sum_{a} X_{u,a} X_{v, \pi_{uv}(a)}$ to be the indicator that the constraint on the edge $(u,v)$ is satisfied.

Then the $Z$ variables satisfy:
\begin{enumerate}
\item Booleanity: $Z_{v,s}^2 = Z_{v,a}$. 
\item Partition constraints: $Z_{v,s}Z_{v,s'} = 0$ for $s \neq s'$, and $\sum_{s} Z_{u,s} = 1$.
\item Crossing edges violate an assignment: $Z_{u,s} Z_{v,s'} Y_{(u,v)} Y'_{(u,v)} = 0$ for every edge $(u,v) \in E$ and $s \neq s'$. 
\end{enumerate}
\end{fact}
\begin{proof}
The first two items are easily verified via direct computation, using properties of the variables $X_{u,a}$.
We prove that the final property holds.
Since our UG instance is affine, we have that for each $i,j \in E$, $\pi_{ij}(a) = a + h_{ij}$ for some $h_{ij} \in \Sigma$.
 Therefore,
 \begin{align*}
 Z_{i,s} Z_{j,t}Y_{(i,j)}Y'_{(i,j)}
 &= \sum_{a,b,c,d \in \Sigma} X_{i,a}X'_{i,a+s} \cdot X_{j,b}X'_{j,b+t} \cdot X_{i,c}X_{j,c + h_{ij}} \cdot X'_{i,d} X'_{j,d+h_{ij}}\\
 &= 0,
 \end{align*}
 where we derive the final equality from the disjointness constraints (i.e. that $X_{i,a}X_{i,b} = 0$ whenever $a \neq b$), as for the above term to be nonzero we require $a = c$, $d = a + s$, $b = d - t + h_{ij} = a + s - t + h_{ij}$, and also $b = c + h_{ij}$, which implies $a + s - t = c$, a contradiction since $t \neq s$.
This establishes the final property.
\end{proof}

\subsubsection{Shift-Symmetry}\label{sec:symm-prelims}
Next, we define the notion of shift-symmetric functions and pseudodistributions and establish some properties that they
satisfies.
\begin{definition}[Shift-Symmetry]\label{def:shift-symm}
We say that a pseudodistribution $\mu$ is shift symmetric if 
for any monomial $\prod\limits_{i=1}^{D} X_{u_i, a_i}$ and $s\in\Sigma$ it holds that
\[
\pE_{\mu}\left[\prod\limits_{i=1}^{D} X_{u_i, a_i}\right]
=
\pE_{\mu}\left[\prod\limits_{i=1}^{D} X_{u_i, a_i+s}\right].
\]
\end{definition}
We say that a polynomial $p(X_{u_1,\sigma_1},\ldots)$ equivalently also thought of as a function $f: \Sigma^n \rightarrow \R$ is shift-symmetric if $f(X) = f(X+s)$ for all $s \in \Sigma$.

Given a pseudodistribution obtained by the \sos relaxation of sum of squares of an 
Affine Unique-Game, we are able to 
transform it into a shift-symmetric pseudodistributions 
with the same value by defining a new pseudodistribution $\mu^{\text{sym}}$ as follows
\[
\widetilde{E}_{\mu^{\text{sym}}}\left[\prod\limits_{i=1}^{D} X_{u_i, a_i}\right] 
= 
\frac{1}{|\Sigma|}\sum\limits_{s\in\Sigma}
\pE_{\mu}\left[\prod\limits_{i=1}^{D} X_{u_i, a_i+s}\right].
\]
It is easy to check that $\pE_{\mu^{\text{sym}}}$ is a valid pseudoexpectation, 
and that it's value is the same; indeed, this follows since if $(a,b)\in \Sigma$ satisfies a constraint in an affine Unique-Games instance, then $(a+s,b+s)$ also satisfy that 
constraint for any $s\in\Sigma$. More generally one can check that $\pE_{\mu^{\text{sym}}}[p(X)] = \pE_{\mu}[p(X)]$ for all shift-symmetric polynomials $p$.
Thus, we will assume henceforth that our pseudoexpectation and the pseudodistribution can be made shift symmetric without losing the value and in general preserving the pseudoexpectation of shift-symmetric polynomials. 



\subsection{Information Theory}
We will use $\mu|_R$ to denote the marginal distribution of a random variable $R \sim \mu$. We use $TV(A,B)$ to denote the total-variation distance between two distributions $A,B$.

\begin{definition}[Mutual Information]
Given a distribution $\mu$ over $(X,Y)$, the mutual information between $X,Y$ is defined as:
\[I_\mu(X;Y) = D_{KL}(\mu~ || ~\mu|_X \times \mu|_Y),\]
where $D_{KL}$ is the Kullback-Leibler divergence. The conditional mutual information between $(X,Y)$ with respect to a random variable $Z$ is defined as:
\[I(X;Y | Z) = \E_{z \sim Z}[I_{\mu | Z = z}(X;Y)].\]
\end{definition}

\begin{lemma}[Pinsker's inequality]
Given any two distributions $D_1,D_2$:
\[TV(D_1, D_2) \leq \sqrt{\frac{1}{2}D_{KL}(D_1 || D_2)}.\]
Using this we get that for a distribution $\mu$ over $(Y,Y')$:
\[TV(\mu, \mu_Y \times \mu_{Y'}) \leq O(\sqrt{I(Y;Y')}).\]
% Using this we get that given any two random variables $X,Y$:
% \[\Cov(X,Y) \leq O(\sqrt{I(X;Y)}).\]
\end{lemma}

\begin{lemma}[Data processing inequality]
Let $X,Y,A,B$ be random variables such that $H(A|X) = 0$ and $H(B|Y) = 0$, i.e. $A$ is fully determined by $X$ and $B$ is fully determined by $Y$. Then:
\[I(A;B) \leq I(X;Y).\]
\end{lemma}







